% Shared packages and formatting-only compatibility definitions.
\usepackage{amsmath,amssymb}
\usepackage{array,booktabs,tabularx}
\usepackage{caption}
\DeclareCaptionFont{bspctablefont}{\fontsize{8}{10}\selectfont}
\captionsetup[table]{font=bspctablefont,labelfont=normalfont,position=top,skip=4pt}
\captionsetup[figure]{position=bottom}
\newcommand{\BSPCtableformat}{\fontsize{8}{10}\selectfont\setlength{\tabcolsep}{2pt}\renewcommand{\arraystretch}{1.12}}
\usepackage{placeins}
\usepackage{flafter}
\usepackage{dblfloatfix}
\usepackage{cite}

% Preserve the approved source behavior: two consecutive references form a range.
\makeatletter
\def\@compress@cite#1#2#3{%
  \ifnum\@cite@incr=\z@
    \advance\@tempcnta\@cclvi
    \ifnum #1=\@tempcnta
      \def\@h@ld{\citedash\@cite@out{#3}}%
      \mathchardef\@cite@incr=\@cclvi
    \else
      \advance\@tempcnta-\@cclv
      \ifnum #1=\@tempcnta
        \def\@h@ld{\citedash\@cite@out{#3}}%
        \mathchardef\@cite@incr=\@ne
      \else
        \@citea\@cite@out{#3}%
      \fi
    \fi
  \else
    \advance\@tempcnta\@cite@incr
    \ifnum #1=\@tempcnta
      \def\@h@ld{\citedash\@cite@out{#3}}%
    \else
      \@h@ld\@citea\@cite@out{#3}%
      \let\@h@ld\@empty
      \mathchardef\@cite@incr=\z@
    \fi
  \fi
  \@tempcnta#1\let\@citea\citepunct
}
\makeatother
\renewcommand{\citedash}{\hbox{--}\penalty\citepunctpenalty}

\usepackage{url}
\usepackage[hidelinks]{hyperref}

\graphicspath{{figures/}}
\newcolumntype{Y}{>{\raggedright\arraybackslash}X}
\newcolumntype{Z}{>{\centering\arraybackslash}X}
\setlength{\emergencystretch}{1em}
\hyphenation{echo-car-di-og-ra-phy}

% The same body file expands to ordinary floats in preprint mode and to
% full-width floats for wide items in the 5p two-column preview. Define the
% selected environment statically so a conditional never crosses a float group.
\ifBSPCtwocolumn
  \newenvironment{bspcwidefigure}[1][!t]{\begin{figure*}[#1]}{\end{figure*}}
  \newenvironment{bspclabelefficiencyfigure}{\begin{figure*}[t]}{\end{figure*}}
  \newenvironment{bspcwidetable}[1][!t]{\begin{table*}[#1]}{\end{table*}}
  \newcommand{\BSPCcompacttableformat}{\BSPCtableformat}
  \newcommand{\BSPCstrictnopretraining}{\shortstack[l]{Strict no-\\pretraining}}
  \newcommand{\BSPCstorelabelefficiencyinterpretation}[1]{\gdef\BSPClabelefficiencyinterpretation{#1}}
  \newcommand{\BSPCemitlabelefficiencyinterpretation}{\BSPClabelefficiencyinterpretation}
  % Keep prose available to backfill pages while wide floats wait for their
  % next legal top position. Source order and the double-float limit retain the
  % architecture sequence without forcing a sparse transition page.
  \newcommand{\BSPCseparatearchitecturefigures}{}
  \newcommand{\BSPCseparatelabelefficiencyfloats}{}
  \newcommand{\BSPCseparateresultfloatgroups}{}
  \newcommand{\BSPCkeependpointorder}{}
\else
  % Preprint mode can use here/top/bottom/float-page placement.  The body keeps
  % the same scientific float declarations as the two-column build.
  \newenvironment{bspcwidefigure}[1][!htbp]{\begin{figure}[!htbp]}{\end{figure}}
  \newenvironment{bspclabelefficiencyfigure}{\begin{figure}[!htbp]}{\end{figure}}
  \newenvironment{bspcwidetable}[1][!htbp]{\begin{table}[!htbp]}{\end{table}}
  \newcommand{\BSPCcompacttableformat}{\BSPCtableformat}
  \newcommand{\BSPCstrictnopretraining}{Strict no-pretraining}
  \newcommand{\BSPCstorelabelefficiencyinterpretation}[1]{#1}
  \newcommand{\BSPCemitlabelefficiencyinterpretation}{}
  \newcommand{\BSPCseparatearchitecturefigures}{}
  \newcommand{\BSPCseparatelabelefficiencyfloats}{}
  \newcommand{\BSPCseparateresultfloatgroups}{}
  \newcommand{\BSPCkeependpointorder}{}
\fi
